\section{Mandatory-service evidence and closure}\label{service-evidence}

The proposed \texttt{service-evidence-v1} mailbox covers protocol-assigned
A detection and B transformation. The two-ticket exercise below validates
transport and review of its B output; its sequential A fallbacks do not
constitute the complete common detector panel. An emission-bearing
benchmark needs a separately frozen schedule for that matrix and its
dependent requests. Private customer jobs retain their paid
lifecycle and never enter mandatory-service counts. Full encrypted payloads
are published on chain; a hash or HTTP locator alone is insufficient.
Publication establishes bytes, sender authentication, and timing. Private
validity and quality require the certificates below. This mailbox is a
specified pilot, not an implemented throughput result.

\paragraph{Frozen authority.}
Before revealing task contents, exclude the exact assigned A/B UIDs from
the round's eligible validator set. Bind every remaining UID, registration
generation, hotkey, signing and encryption key versions, and integer native
effective consensus stake, including native alpha/TAO scaling, to one
finalized state snapshot. Its manifest names the block hash, runtime and
storage schema, weight field and units, complete records, and total weight
\(W\). Initial loading and later snapshots require verified finalized state
proofs or a chain-verified native snapshot. A relayer's assertion alone is
insufficient~\cite{btprecompilelifecycle}.

Every final validity, quality, or incident verdict requires signatures
whose distinct frozen weights sum to \(w\), with \(3w>2W\). Assume
dishonest weight is strictly below \(W/3\) in this conditional set.
Missing keys, nonresponse, recusal, and later key changes never shrink
\(W\). Three hash-assigned preparers organize evidence; they cannot
authorize a verdict. Honest validators independently check evidence and
sign at most one verdict per stage. Conflicting strict quorums would share
more than \(W/3\) weight; a detected conflict halts affected settlement.
Signatures bind the verdict itself, evidence, snapshot, stage, and scope.

\paragraph{Fixed pilot.}
The instance manifest supplies actual chain/contract addresses, verified
snapshots, authorizations, service/schema hashes, and funded reservations.
Missing bindings, zero eligible weight, or insufficient qualified key and
verification capacity prevent issue. At native epoch start \(H\), issue at most
two tickets, one A and one B, in a 360-block measurement epoch. Each has
three prescribed providers total. Freeze these settings before issue:

\begin{center}
\begin{tabular}{@{}ll@{}}
\toprule
Parameter & Value \\
\midrule
Input / output cap & 4,096 / 8,192 bytes \\
Ciphertext chunk / publication cap & 4 KiB / 8 MiB per stage \\
Global epoch publication cap & 64 MiB, including all witness copies \\
Request publication / certificate & \(H+12\) / \(H+24\) \\
Provider A / B response window & 12 / 36 blocks \\
Certificate window per attempt & 12 blocks \\
Preparer commitment / private opening & \(H+180\) / \(H+192\) \\
Final verdict and incident close & \(H+204\) \\
\bottomrule
\end{tabular}
\end{center}

For response window \(r\), attempt \(j\in\{1,2,3\}\) starts at
\(H+24+(j-1)(r+12)\), publishes by start plus \(r\), and certifies by
start plus \(r+12\). Final A/B attempt cutoffs are \(H+96\)/\(H+168\).
Early declines do not advance later windows. Reserve all possible bytes,
gas, response and verification capacity; insufficiency prevents issue and
does not authorize another draw. Two tickets provide a bounded transport
pilot, not availability evidence for every UID. Settlement names a later
native payout cycle and its actual commit/reveal configuration; no claim
is made that reveal delays fit within the unused part of this epoch.

\paragraph{DATA, then witness.}
Use HPKE base mode with X25519/HKDF-SHA256/ChaCha20Poly1305, fresh
single-shot ephemeral material per recipient, and separate registered
sender signatures~\cite{hpke}. Deterministic CBOR records reject duplicate
keys, indefinite encodings, tags, and floats~\cite{cbor}. A salted payload
commitment and encryption context bind protocol, chain, contract, ticket,
stage, attempt, both UID generations, and recipient key version. Input
DATA copies address the requester and three prescribed providers; response
copies address the requester and actual provider. Curator authorization
covers those recipients and the complete validator set. Hidden labels and
review answers remain in validator-only evidence packets.

First fix and publish the signed DATA manifest and all ciphertext chunks.
Then form a witness containing that DATA digest, plaintext, salt, and each
DATA envelope's ephemeral generation material. Encrypt the witness
separately to every frozen validator and publish those complete ciphertexts
under a second manifest referencing DATA. Validators reconstruct each DATA
encryption and compare the exact bytes and recipient bindings, as well as
authorization, commitment, and schema. Decrypting only their own copy would
not check the miner's copy. DATA never references a witness hash; witnesses
do not contain their own encryption secrets. This avoids a hash cycle.
Generation witnesses remain confidential and need implementation/security
review. Their certificates do not assert that every validator decrypted
its witness copy. Public ciphertext remains exposed to future key
compromise; authorized archives retain complete replay evidence.

\paragraph{Transitions and attribution.}
Issue permanently consumes the requester index and activates its
obligation. One immutable manifest fixes each stage; identical retries
are idempotent and conflicting signed manifests remain evidence. A
requester's decline, absent complete publication, or certified invalid
input is its failure. An authenticated decline included by the stage's
publication cutoff is terminal for that stage. Retain later bytes as
evidence, but reject an acceptance certificate that ignores the decline;
late declines cannot reopen closed accounting. Only a timely global input-acceptance certificate
activates the first provider at its fixed start. No input acceptance
means no provider obligation. A provider's decline, absent publication,
or globally certified invalid response permits the next fixed fallback;
successful fallback never erases preceding failures. Unresolved validity
stops dependent work without another draw. Late answers remain late.

Complete publication without a global validity verdict closes unresolved;
recipient accusation alone cannot assign fault. Keys must be qualified
before the snapshot. Key loss afterward cannot blame a sender whose
encryption was certified correct; rotations and UID reuse cannot change
issued tickets. A timely dishonest answer may satisfy service availability
while failing quality. Insufficient final review evidence voids the shared
quality comparison consistently, without deleting attributable service
failures. The detailed record and certificate rules are specified in the
companion service-evidence document.

\paragraph{Counts and finality.}
For each mandatory interface/role, retain activated obligations
\(A=S+F+U+V\): success, attributable failure, unresolved, and platform void.
Unused reservations are separate. Let \(N=S+F\); require \(U=0\), \(N>0\),
and \(10F\leq N\) for the class and aggregate availability gates. Publish
all counts. Keep \(D_e=\max(0,D_{e-1}+10F-N)\); neither unresolved nor void
work retires a deficit. Require zero outstanding deficits and current
qualification; customer jobs and other classes cannot dilute these gates.

Only canonical finalized inclusion counts, with deadline equality timely.
Subtensor separates block production from GRANDPA
finality~\cite{btchainconsensus}. A local timeout or censorship allegation
does not prove a platform exception. Normal operation assumes fair
inclusion and sufficiently prompt finality. A finality stall delays
observing closure without extending block-height deadlines. A strict
global incident certificate must be included by \(H+204\), even if its
block finalizes later, and identify the affected interval and every
intersecting ticket/comparison. The rule
voids all activated obligations in that scope, including successes and
failures, and cannot selectively pardon a miner, shrink \(W\), or reopen
closed epochs. Without that certificate, normal inclusion rules apply;
missing verdicts remain unresolved. Gas, history retrieval, proof-verified
snapshots, and the declared native settlement timing must be demonstrated
before this EVM pilot issues work~\cite{evm}.
